\documentclass[11pt,a4paper]{article}
\usepackage[margin=2.2cm]{geometry}
\usepackage[T1]{fontenc}
\usepackage{amsmath,amssymb,xcolor,enumitem,booktabs,graphicx,tabularx}
\usepackage[normalem]{ulem}
\setlength{\parskip}{0.6em}\setlength{\parindent}{0pt}
\definecolor{newc}{HTML}{1f6f43}
\newcommand{\new}[1]{\textcolor{newc}{#1}}
\newcommand{\old}[1]{\textcolor{gray}{\sout{#1}}}
\newcommand{\note}[1]{{\small\itshape [#1]}}
\newcommand{\lab}[1]{\textsf{\footnotesize\detokenize{#1}}}
\renewcommand{\ref}[1]{\lab{#1}}
\renewcommand{\eqref}[1]{(\lab{#1})}
\newcommand{\nref}[1]{\lab{#1}}
\newcommand{\citep}[1]{[\lab{#1}]}
\newcommand{\citet}[1]{\lab{#1}}
\newcommand{\mockup}[1]{\fbox{\small mockup: #1}}
\title{Proposed change, Method parts that belong elsewhere}
\date{25 November 2026, version 1}
\begin{document}\sloppy
\maketitle

\section*{How to read this}
The Token Passing review is applied as [A137]. This review moves the other
parts of the Method and the Implementation chapter that are not method. Each
part either is a result, a value, or experiment design, or it is said twice.
Nothing in the model changes.

\begin{enumerate}[nosep]
  \item The lock example from an earlier run moves to the Results.
  \item The fleet 32 story in the reward moves to Appendix~I.
  \item Training: design stays in the Method, seeds and step counts go to
    the Experiments.
  \item ``What is logged'': the two sentences on scoring values go.
  \item ``Computational cost'': why and how runtime is reported goes to the
    Experiments.
  \item Assignment: the put-away guess goes, the Discussion already has it.
  \item Small: ``never overridden'' said once less, the Ma et al.\ match in
    numbers.
\end{enumerate}

Black text stays, \new{green} is new and \old{struck} text goes.
\note{Notes in brackets} say why. Nothing below is in the thesis yet.

% ------------------------------------------------------------------
\section*{1. The lock example moves to the Results}

\textbf{Method, \lab{sec:method:resolution}.} The paragraph on how a lock
happens in a one-lane aisle and the paragraph on why no rule breaks it stay
as they are. The figure and the paragraph that tells the earlier run go.

\old{[Figure \lab{fig:workstationlock}: A lock at workstation $W_1$ with eight
agents, under the decentralised controller without messages in an earlier
run.]}

\old{Figure~\ref{fig:workstationlock} shows such a lock in a test warehouse
with eight agents. All four workstations of one category hang off the outer
aisle, \ldots\ so the whole fleet finishes only two more tasks in the last
843 steps of the episode.} \note{An observed run is a result. The Method keeps
the mechanism.}

\textbf{Method, Token Passing.} \ldots so Token Passing never meets the lock
at a workstation \old{of Figure~\ref{fig:workstationlock}}\new{described in
Section~\ref{sec:method:resolution}}. \note{The figure is now in the Results,
so the Method points back, not forward.}

\textbf{Experiments, Results, after ``Token Passing never locks.''}

\new{Figure~\ref{fig:workstationlock} shows such a lock in an earlier run, in a
test warehouse with eight agents.} All four workstations of one category hang
off the outer aisle, which is one lane wide. At $t = 900$ the work still
flows, \ldots\ so the whole fleet finishes only two more tasks in the last
843 steps of the episode. \note{Text and figure unchanged, only the first
sentence. The lock definition just above changes with [A136], so this move
goes in together with it.}

% ------------------------------------------------------------------
\section*{2. The fleet 32 story moves to Appendix I}

\textbf{Method, \lab{sec:method:rl}.}

The reason is that an agent avoids a charge on finishing by never finishing.
\old{When a new task's distance was charged on the step it was handed out,
delivering cost $30$ to $60$ at once against a bonus of $10$, and agents in a
fleet of $32$ learned to stop just before the workstation.}
\new{If a new task's distance were charged on the step it is handed out,
delivering would cost about the length of the next task at once, more than
the delivery bonus (Appendix~\ref{app:reward:method}).} \note{The reason stays
in words. The run that showed it is a result.}

\textbf{Appendix I, \lab{app:reward:method}, a new paragraph before the step
costs.}

\new{The progress term restarts with every task. An earlier version charged
a new task's distance on the step it was handed out. Delivering then cost
$30$ to $60$ at once against a bonus of $10$, and in training runs at $32$
agents the agents learned to stop just before the workstation. Counting each
task from its own start removed that charge
(Section~\ref{sec:method:rl}).}

% ------------------------------------------------------------------
\section*{3. Training: design in the Method, numbers in the Experiments}

Both places now say the PPO budget, one policy per storage rule and the
seeds. The Method keeps the choices and their reasons. The Experiments keeps
the seeds, the step counts and how a checkpoint is chosen.

\textbf{Method, \lab{sec:method:training}, second paragraph.}

\old{Each seed is its own warehouse. Before any run starts the seeds are split
three ways. Training seeds are played in training, validation seeds choose the
checkpoint and test seeds give the results. Training plays a new warehouse in
every episode, so a policy cannot learn one map by heart.}
\new{Each seed is its own warehouse, and training plays a new one in every
episode, so a policy cannot learn one map by heart.}
\old{At 16 agents, the fleet of the main results, training}\new{Training} has two
stages, the same for every controller and storage rule. \old{A first run plays
500\,000 steps on episodes of 2000. A second run starts from the weights of
the first and plays 2\,000\,000 steps on episodes of 16\,000. Starting there
saves training time. The episodes of the second run are longer than an
evaluation run, 10\,000 steps, so}\new{A first run plays short episodes. A
second run starts from its weights, which saves training time, and plays
episodes longer than a test run. So} a fleet that falls behind over \old{an
evaluation}\new{a test} run also falls behind in training and in the
validation that chooses the checkpoint. \new{The step counts are in
Section~\ref{sec:exp:setup}, and} how the two stages were chosen is in
Appendix~\ref{app:training}. \note{The reasons stay, the numbers go.}

\textbf{Experiments, ``Training''.}

\old{Each learned controller is trained with PPO under the same budget of
environment steps, using the values in Table~\ref{tab:trainparams}. Each
storage rule has its own policy, trained under the rule it is tested with
(Section~\ref{sec:method:training}).}
\new{The PPO values and the budget of environment steps are in
Table~\ref{tab:trainparams}. At 16 agents, the fleet of the main results, the
first stage of training plays 500\,000 steps on episodes of 2000 and the
second 2\,000\,000 steps on episodes of 16\,000, against a test run of
10\,000 steps. The seeds are split three ways before any run starts. Training
seeds are played in training, validation seeds choose the checkpoint and test
seeds give the results.} \old{Training seeds and evaluation seeds are counted
apart.} The spread over training seeds shows \ldots\ [rest unchanged: the
mockups and how a checkpoint is chosen]. \note{The PPO method and one policy
per rule are said in the Method. The split of seeds sits beside the
checkpoint rule that uses it.}

% ------------------------------------------------------------------
\section*{4. ``What is logged''}

\ldots Appendix~\ref{app:repro} lists what is logged and which measure each
record gives. \old{Table~\ref{tab:evalparams} lists the values used to score a
run. The load levels and $f_{\mathrm{up}}$ come from a pilot of Token Passing
(Appendix~\ref{app:cal:load}).} \note{Scoring values are experiment design.
``Scoring a condition'' already cites the table, and ``Factors'' now says the
load comes from Token Passing.}

\textbf{Experiments, ``Scoring a condition''.} \ldots The keep-up check of
\nref{eq:throughput} comes first\new{, with the share $f_{\mathrm{up}}$ set in
Appendix~\ref{app:cal:load}}. \ldots \note{Where $f_{\mathrm{up}}$ comes from
was said only in the sentence that goes.}

% ------------------------------------------------------------------
\section*{5. ``Computational cost''}

\textbf{Implementation, \lab{sec:impl:cost}.}

\old{Runtime is one of the measured results.} \new{Runtime $\kappa_T$ is timed
in the simulator.} The time per decision is logged separately for
\nref{eq:assignmentrule} and for routing, on the same machine for every
controller. The learned controllers run their policy on \old{its}\new{the} GPU.
Token Passing runs on the CPU, the hardware it is built for, so each method is
timed on its own hardware. \old{For the learned controllers, training time and
decision time are reported separately, because the practical case for a
learned router rests on fast decisions, not on fast training. Token Passing
needs no training, and its planning time per decision is what the learned
routers' decision time is compared with.} \note{How runtime is timed stays.
What is reported and why moves.}

\textbf{Experiments, ``Scoring a condition'', after ``\ldots crowding, the
mean service time of put-aways and runtime are reported next to it.''}

\new{Runtime is reported because a learned controller is only of practical use
if it decides fast enough. So for the learned controllers training time and
runtime are reported apart. Token Passing needs no training, and its runtime
is what theirs is compared with.} \note{One name, ``runtime'', instead of
decision time and planning time (implementation comment C1). ``Learned
controllers'', not ``learned routers''.}

\textbf{Experiments, ``Compute environment''.}

All training and evaluation ran on the machine in
Table~\ref{tab:exp:compute}. Runtime $\kappa_T$ is measured on \old{the same
machine for every controller, the learned controllers with their policy on the
GPU and Token Passing on the CPU}\new{this machine}
(Section~\ref{sec:impl:cost}). \note{The hardware was told in both places
(implementation comment D2).}

% ------------------------------------------------------------------
\section*{6. Assignment, the put-away guess}

\textbf{Method, \lab{sec:method:assignment}.}

Picks and put-aways wait in the same queue and are taken alike. A put-away
gets no priority of its own, because oldest first is the simplest rule and
adds no parameter. \old{A put-away that waits keeps its units from being
picked, so at a high load a priority for put-aways might shorten the wait for
stock. Section~\ref{sec:exp:discussion} returns to this.} Tasks released in
the same timestep go in order of task index. \ldots \note{The Discussion item
``Priority for put-aways'' already says this in full. The Method keeps the
reason for the choice.}

% ------------------------------------------------------------------
\section*{7. Small}

\textbf{``Never overridden'' said once less.} It is in Choosing together
(``Conflict resolution therefore never overrides a centralised agent''), in
When moves still clash, and in the reward (who pays the override term). The
reward sentence stays, because it says who pays. The second goes.

When moves still clash, last sentences: The claim mask of the centralised and
section-based controllers is stricter, because an agent choosing there does
not yet know whether a later agent will leave. \old{So a centralised joint
action is always accepted whole, and a section-based one is changed only at
the border between two zones.} \note{Said in Choosing together, both for the
centralised agent and at the border.}

\textbf{The Ma et al.\ match.} Correctness and reproducibility: Token Passing
in this simulator also \old{reproduces the published results of
\citet{ma2017lifelong_mapf} to within about 10\%}\new{gives service times
close to those published by \citet{ma2017lifelong_mapf}, $8\%$ shorter on
average} (Appendix~\ref{app:repro:ma}). Appendix~K, first sentence of
\lab{app:repro:ma}, the same change. \note{Appendix~K gives $8\%$ shorter on
average, from $34\%$ shorter to $3\%$ longer. ``Within about $10\%$'' reads as
every setting (implementation comment A7).}

% ------------------------------------------------------------------
\section*{Not in this review}
The rest of \texttt{implementation\_comments.md} (A2, A3, A5, A6, A8, B1, C2,
E1) changes wording inside the Implementation chapter, not where text sits.
A4 there (the load levels in \lab{tab:evalparams} against Appendix~E) is a
question about values.

\section*{Questions}
\begin{enumerate}[nosep]
  \item Part 1. With the figure gone, the Method explains the lock in words
    only. Should it get a small drawing in the viewer style (one-lane aisle,
    $a_1$ on the workstation, $a_2$ in front, two agents behind), or are the
    words enough?
  \item Part 1 waits for the new lock definition ([A136]). Apply parts 2 to 7
    now and part 1 with [A136], or all together?
  \item Should the rest of the implementation comments (see above) come in a
    second version of this review?
\end{enumerate}

\end{document}
